%% ============================================================
%%  SECCIÓN 8 — Análisis Experimental
%% ============================================================

\section{Análisis Experimental}
\label{sec:analisis_exp}

El análisis cualitativo de las observaciones registradas durante la experiencia permite identificar el comportamiento de las magnitudes físicas y la influencia de las condiciones instrumentales sobre la fidelidad de las medidas:

\subsection{Efecto de la Sensibilidad Vertical en la Medición de Tensiones DC}
Al evaluar las lecturas de voltaje continuo en la retícula del osciloscopio (Tabla~\ref{tab:resultados_dc}), se evidencia que la exactitud visual está fuertemente condicionada por la escala de voltios por división seleccionada:
\begin{itemize}
    \item \textbf{Escalas no óptimas (\qty{5}{\volt\per\division}):} Para señales de bajo potencial (como la pila de \qty{1,5}{\volt}), la deflexión vertical ocupa apenas $0{,}3\,\text{divisiones}$. En esta condición, una pequeña incertidumbre de apreciación ocular ($\pm 0{,}1\,\text{div}$) representa una variación relativa considerable ($\pm\qty{0,5}{\volt}$ o un \qty{33}{\percent} de la magnitud medida), incrementando la dispersión aparente de los datos.
    \item \textbf{Escalas óptimas (\qty{0,2}{\volt\per\division} y \qty{0,5}{\volt\per\division}):} Al incrementar la ganancia vertical de modo que el trazo ocupe entre el \qty{60}{\percent} y el \qty{90}{\percent} de la altura disponible en la pantalla sin saturarse, la misma incertidumbre de $\pm 0{,}1\,\text{div}$ se traduce en un error absoluto mínimo ($\pm\qty{0,02}{\volt}$ a $\pm\qty{0,05}{\volt}$), logrando una coincidencia visual casi perfecta con la lectura del multímetro de referencia.
    \item \textbf{Saturación y fuera de escala:} Cuando la tensión de entrada supera el rango dinámico vertical disponible ($> 8\,\text{divisiones}$ útiles), el haz electrónico digitalizado sobrepasa los límites de la cuadrícula, provocando recorte de señal. Por ello, el operador debe seleccionar la menor escala vertical posible que mantenga la señal íntegramente visible.
    \item \textbf{Comportamiento de la fuente en circuito abierto (en vacío):} Las tensiones medidas en las salidas fijas de la fuente continua resultaron sistemáticamente superiores a sus valores nominales de placa (\qty{5,220}{\volt} frente a \qty{4,50}{\volt}, \qty{8,030}{\volt} frente a \qty{7,50}{\volt} y \qty{11,150}{\volt} frente a \qty{10,50}{\volt}). Una explicación plausible es el comportamiento característico de fuentes de banco lineales en vacío (filtrado capacitivo sin caída óhmica de carga, $V_{\text{out}} \approx V_{\text{th}}$); sin embargo, dado que no se realizaron mediciones bajo distintas cargas resistivas durante la práctica, esto se formula como hipótesis tentativa y no como causa demostrada. Por tanto, la referencia rigurosa de contraste para la calibración del osciloscopio es el voltaje real registrado simultáneamente por el multímetro patrón ($V_{\text{mult}}$).
\end{itemize}

\subsection{Comportamiento de la Tensión Alterna del Transformador}
La señal proveniente del transformador reductor presenta una morfología estrictamente sinusoidal y simétrica respecto al eje neutro, lo que confirma la ausencia de componentes continuas significativas ($\text{DC} \approx \qty{0}{\volt}$).

En el dominio temporal, al reducir la base de tiempo de \qty{5}{\milli\second\per\division} a \qty{2}{\milli\second\per\division}, la onda se expande horizontalmente abarcando $8{,}3\,\text{divisiones}$ por ciclo en lugar de $3{,}3\,\text{divisiones}$. Esta expansión reduce sustancialmente el error de apreciación en la determinación de los puntos de cruce por cero y la cúspide de la onda, mejorando la precisión en el cálculo del período fundamental $T$.

\subsection{Morfología de Señales Sintéticas del Generador OWON}
Las tres funciones periódicas evaluadas a \qty{250}{\hertz} demostraron alta estabilidad temporal gracias al principio de síntesis digital directa (DDS) del generador:
\begin{itemize}
    \item La \textbf{onda senoidal} no exhibe distorsiones apreciables visualmente en la escala empleada.
    \item La \textbf{onda triangular} presenta pendientes ascendentes y descendentes de igual magnitud absoluta, con inversión de signo abrupta en los vértices pico.
    \item La \textbf{onda cuadrada} muestra transiciones de flanco de conmutación sumamente pronunciadas (tiempo de subida nominal $t_r < \qty{15}{\nano\second}$ especificado por el fabricante~\cite{owon_ags_spec}), manteniendo niveles de tensión estacionarios de $+V_p$ y $-V_p$ durante semiperíodos simétricos (ciclo de trabajo del \qty{50}{\percent}).
\end{itemize}

\subsection{Dinámica Geométrica en el Modo Graficador X-Y}
En la modalidad $X$-$Y$, la combinación ortogonal de señales senoidales revela directamente la correlación de frecuencias y el desfase temporal:
\begin{itemize}
    \item Para la relación $1:1$ (\qty{60}{\hertz} : \qty{60}{\hertz}), la curva genera una elipse estacionaria. Al aplicar la función $\text{CH2 INVERT}$, la deflexión en el eje $Y$ se multiplica por $-1$, produciendo una reflexión especular respecto al eje horizontal $X$ que invierte la inclinación de la elipse.
    \item Al elevar la frecuencia del canal vertical a múltiplos enteros exactos ($120$, $180$ y \qty{240}{\hertz}), la trayectoria se pliega sobre sí misma originando figuras de Lissajous con $2$, $3$ y $4$ lóbulos verticales cerrados y estables.
    \item Para tensiones continuas constantes en ambos canales, al no existir componente armónica dependiente del tiempo, el barrido se reduce a un punto fijo estático en el plano cartesiano $(V_x, V_y)$, el cual migra al cuarto cuadrante $(V_x, -V_y)$ al invertir el canal $Y$.
\end{itemize}
